\begin{enonce}{Lemma}{}
The $S$-derivations of $u$ are exactly the $S$-deviations of $u$ of order $1$ which vanish on the unit section of $\OX$; they correspond to the $\OS(S)$-module
\[
\Hom_{\OS}(\mathscr I_u/\mathscr I_u^2,\OS),
\]
and there is an isomorphism of $\OS(S)$-modules $\operatorname{D\acute ev}^{\leq1}(u)\simeq\OS(S)\oplus\operatorname{D\acute er}_S(u)$.
\end{enonce}
Returning to the general case, one deduces from this, with the notation of 1.3,
\begin{enonce}{Corollary}{}
Let $u:Y\to X$ be an $S$-morphism and $\Gamma u:Y\to Y\times X$ its graph. There is a canonical isomorphism of $\OY(Y)$-modules
\[
\operatorname{D\acute er}_S(u)\simeq\operatorname{D\acute er}_Y(\Gamma u)\simeq\Hom_{\OY}(\mathscr I_{\Gamma u}/\mathscr I_{\Gamma u}^2,\OY).
\]
\end{enonce}

\begin{definition}{1.4}\label{VIIA.1.4}
Let $X$ be an $S$-scheme. An \emph{$S$-differential operator} (resp.\ an $S$-differential operator of order $\leq n$) on $X$ means any $S$-deviation (resp.\ any $S$-deviation of order $\leq n$) of the identity morphism of $X$.

According to 1.1, an $S$-differential operator of order $\leq n$ is thus an endomorphism of the $p^{-1}(\OS)$-module $\OX$ satisfying equalities $(*_n)$ of 1.1. We shall denote by $\mathrm{Dif}_{X/S}^n$ the $\Gamma(\OS)$-module\nde{In this expos\'e, the ring $\Gamma(S,\OS)=\OS(S)$ is denoted by $\Gamma(\OS)$.} consisting of the $S$-differential operators of order $\leq n$, and by $\mathrm{Dif}_{X/S}$ that consisting of all the $S$-differential operators.

As we saw in 1.2, one can compose the $S$-deviations of $\mathrm{id}_X$, which equips $\mathrm{Dif}_{X/S}$ with a $\Gamma(\OS)$-algebra structure; we shall say that it is the \emph{algebra of differential operators of $X/S$}.

Similarly, for every open subset $V$ of $X$, put $\mathscr D\!\mathrm{if}_{X/S}(V)=\mathrm{Dif}_{V/S}=\operatorname{D\acute ev}(\mathrm{id}_V)$; according to 1.1.3, this defines a sheaf of $\OX$-modules, called the \emph{sheaf of $S$-differential operators on $X$}.\nde{Here the original has been modified; it mentioned the sheaf $U\mapsto\mathrm{Dif}_{X_U/U}$, where $U$ ranges over the open subsets of $S$; this is the direct image of $\mathscr D\!\mathrm{if}_{X/S}$ under the morphism $p_X:X\to S$.}
\end{definition}

\numpara{1.4.1}
As we saw in 1.3, one can interpret the differential operators of $X/S$ by means of the graph of the identity morphism of $X$, that is, of the diagonal morphism $\Delta=\Delta_{X/S}$ from $X$ to $X\times X$. Let us translate the statements of 1.3 into the present context.

Equip $\OO_{X\times X}$ with the $\mathrm{pr}_1^{-1}(\OX)$-algebra structure defined by $\mathrm{pr}_1$, so that $\Delta^{-1}(\OO_{X\times X})$ is equipped with an algebra structure over $\OX=\Delta^{-1}\mathrm{pr}_1^{-1}(\OX)$. Let $\mathscr I_{X/S}$ be the kernel of the homomorphism
\[
\Delta^{-1}(\OO_{X\times X})\longrightarrow\OX
\]
adjoint to the homomorphism $\OO_{X\times X}\to\Delta_*(\OX)$, and let $\mathscr P_{X/S}^m$ be the $\OX$-algebra
\[
\Delta^{-1}(\OO_{X\times X})/\mathscr I_{X/S}^{m+1}.
\]
If $V$ is an affine open subset of $S$ and $U$ an affine open subset of $X$ over $V$, and if one puts $k=\Gamma(V,\OS)$ and $A=\Gamma(U,\OX)$, one therefore has:
\[
\Gamma(U,\mathscr P_{X/S}^m)=(A\otimes_k A)/I^{m+1},
\]
where $I$ is the ideal generated by the elements $a\otimes1-1\otimes a$, for $a\in A$. This being so, according to 1.3 there is \emph{an isomorphism of $\OX(X)$-modules}:
\[
j_X:\mathrm{Dif}_{X/S}^m\isomto\Hom_{\OX}(\mathscr P_{X/S}^m,\OX)
\]
which may be defined as follows: if $d$ belongs to $\mathrm{Dif}_{X/S}^m$ and $c$ is a section of $\mathscr P_{X/S}^m$ over $U$ of the form $a\otimes b+I^{m+1}$, one has $j_X(d)(c)=a\cdot d(b)$.\nde{Via this isomorphism, the $X$-derivations of $\Delta_{X/S}$ correspond, according to 1.3.1, to the $S$-derivations of $\mathrm{id}_X$, \ie to the $p^{-1}(\OS)$-derivations of $\OX$.}

\numpara{1.4.2}
Let $d$ be a differential operator and $u$ a section of $X$ over $S$. We call the composite $S$-deviation
\[
\xymatrix{S\ar[r]^u&X\ar[r]^d_{\mathrm{id}_X}&X.}
\]
the \emph{value of $d$ at $u$}.

According to 1.3 and 1.4.1, if $d$ is a differential operator of order $\leq m$, then $du$ (resp.\ $d$) is canonically associated to a morphism of $\OS$-modules $d':u^{-1}(\OX)/\mathscr I_u^{m+1}\to\OS$ (resp.\ a morphism of $\OX$-modules $d'':\mathscr P_{X/S}^m\to\OX$).

It is clear that one can construct $d'$ from $d''$ in the following way: the square
\[
\xymatrix@C=4em{X\simeq S\times X\ar[r]^{u\times X}\ar[d]_p&X\times X\ar[d]^{\mathrm{pr}_1}\\
S\ar[r]_u&X}
\]
is cartesian, which allows one to identify $X$ with $S\times_X(X\times X)$, $u$ with $S\times_X\Delta$, and hence $u^*(\mathscr P_{X/S}^m)$ with $u^{-1}(\OX)/\mathscr I_u^{m+1}$. One thus identifies $u^*(d'')$ with a morphism
\[
u^{-1}(\OX)/\mathscr I_u^{m+1}\longrightarrow\OS,
\]
which is none other than $d'$.

\numpara{1.5}
Put, as usual, $I_S=\Spec\OS[T]/(T^2)$. Let $\tau:S\to I_S$ be the zero section and $\sigma$ the canonical deviation of $\tau$ which we defined in 1.2.0, \ie the homomorphism of $\OS$-modules which vanishes on the unit section of $\OS[T]/(T^2)$ and sends the class $t$ of $T$ modulo $T^2$ to the unit section of $\OS$.

Let $X$ be an $S$-scheme. To every $I_S$-automorphism $u$ of $I_S\times X$ inducing the identity on $X$, a differential operator $D_u$ of $X$ is associated by composition:
\[
X\simeq S\times X\xrightarrow{\sigma\times X}I_S\times X\xrightarrow{u}I_S\times X\xrightarrow{\mathrm{pr}_2}X.
\]
According to II, 3.14, \emph{the map $u\mapsto D_u$ is an isomorphism of the $\Gamma(\OS)$-Lie algebra}
\[
\Lie(\underline{\Aut}\,X):=\uLie(\underline{\Aut}\,X)(S)
\]
\emph{onto the $\Gamma(\OS)$-Lie algebra of $p^{-1}(\OS)$-derivations of $\OX$}. The inverse isomorphism associates to every derivation $D$ the automorphism of $I_S\times X$ corresponding to the automorphism $a+bt\mapsto a+(Da+b)t$ of $\OX[T]/(T^2)$.

\sgasection{2}{Invariant differential operators on group schemes}
\numpara{2.1}
Let $G$ be an $S$-group scheme; we denote by $\varepsilon$ or $\varepsilon_G:S\to G$ the unit section of $G$.

\begin{enoncerm}{Definition}{}
Let $U(G)$ be the $\Gamma(\OS)$-module of $S$-deviations of $\varepsilon_G$ (or \emph{$S$-deviations of the origin}) (cf.\ 1.1).

If $d$ and $e$ are two elements of $U(G)$, $d\times e$ is an $S$-deviation of $\varepsilon\times\varepsilon:S\simeq S\times S\to G\times G$. The image of $d\times e$ under the multiplication morphism $m:G\times G\to G$ (cf.\ 1.2) will be called the \emph{product of $d$ and $e$} and will be denoted by $d\cdot e$.

The $\Gamma(\OS)$-module $U(G)$ is thus equipped with an associative $\Gamma(\OS)$-algebra structure having $\varepsilon_G$ as its unit element (1.1). We shall say that $U(G)$ is the \emph{infinitesimal algebra of $G$}.\nde{One now says ``the algebra of distributions'' (at the origin) of $G$, cf.\ [DG70], \S~II.4, 6.1 and [Ja03], I~7.7.}

As $T$ ranges over the schemes over $S$, the infinitesimal algebra $U(G_T)$ of the $T$-group $G\times T$ clearly varies contravariantly in $T$, so that we shall be able to speak of the \emph{infinitesimal algebra functor}.

As $T$ ranges over the open subsets of $S$, one therefore obtains a presheaf $T\mapsto U(G_T)$ of $\OS$-algebras; moreover, according to 1.1.3, this is a \emph{sheaf}. We shall denote it by $\mathscr U(G)$ and call it the \emph{sheaf of infinitesimal algebras of $G$}.
\end{enoncerm}

The algebra $U(G)$ is also a covariant functor in $G$. Indeed, if $u:G\to H$ is a homomorphism of $S$-groups and $d$ an $S$-deviation of $\varepsilon_G$, the image of $d$ under $u$ is an element $U(u)(d)=ud$ of $U(H)$. The map $U(u):U(G)\to U(H)$ thus defined is clearly a homomorphism of $\Gamma(\OS)$-algebras. One similarly defines a homomorphism $\mathscr U(u)$ from $\mathscr U(G)$ to $\mathscr U(H)$.

\numpara{2.2}
Let $d$ be an element of $U(G)$, \ie an $S$-deviation of the origin of $G$. Consider the $S$-deviation $d\times G$ of $\varepsilon\times G:G\simeq S\times G\to G\times G$ obtained from $d$ by base change (1.2.2); the image of $d\times G$ under the multiplication morphism $m:G\times G\to G$ is an $S$-deviation of $m\circ(\varepsilon\times\mathrm{id}_G)=\mathrm{id}_G$, \ie an element of $\mathrm{Dif}_{G/S}$, which will be denoted by $d^G$.

The map $d\mapsto d^G$ is clearly $\Gamma(\OS)$-linear, and the ``commutative'' diagram below shows that $(e\cdot d)^G=d^G\cdot e^G$:\nde{The original has been corrected by replacing $d\times G\times G$ in the diagram by $G\times d\times G$, so that the composite on the left side of the triangle is $(e\times d)\times G$, and the map $d\mapsto d^G$ is an \emph{anti-isomorphism} from $U(G)$ onto the \emph{right-invariant} differential operators (cf.\ 2.3, 2.4 below); on the other hand, defining ${}^Gd$ as the image under $m$ of $G\times d$, one would similarly obtain an isomorphism from $U(G)$ onto the \emph{left-invariant} differential operators (cf.\ [DG70], \S~II.4, Th.\ 6.5). Sections 2.4 and 2.5 have been corrected accordingly.}
\[
\xymatrix@C=1.1em@R=2.2em{
&&G\times G\times G\ar[rr]^{m\times G}\ar[dr]^{G\times m}&&G\times G\ar[dd]^m\\
&G\times G\ar[ur]^{G\times d\times G}_{\varepsilon\times G\times G}\ar[dr]^m&&G\times G\ar[dr]^m&\\
G\ar[ur]^{e\times G}_{\varepsilon\times G}\ar[rr]^{e^G}_{\mathrm{id}_G}&&G\ar[ur]^{d\times G}_{\varepsilon\times G}\ar[rr]^{d^G}_{\mathrm{id}_G}&&G.}
\]
The commutativity of the two lower triangles follows indeed from the definition of $d^G$ and $e^G$; on the other hand, the composite $S$-deviation of $e\times G$ and $G\times d\times G$ is $(e\times d)\times G$ (cf.\ 1.2.2), its image under $m\times G$ is $(e\cdot d)\times G$, and the image of the latter under $m$ is therefore equal to $(e\cdot d)^G$.

One thus obtains an anti-homomorphism $U(G)\to\mathrm{Dif}_{G/S}$ of $\Gamma(\OS)$-algebras, called \emph{right translation}.\nde{It would be preferable to call it \emph{left operation}. Indeed, let $d$ be, for example, an $S$-derivation of the origin; according to 1.2.1, $d$ is the composite of the $S$-derivation $(\tau,\partial_t):S\to I_S$ and a morphism $x:I_S\to G$ such that $x\circ\tau=\varepsilon$ (\ie $x\in\uLie(G/S)(S)$), and then $d^G$ is the derivation of $\OO_G$ sending a local section $\phi$ to the section $g\mapsto\partial_t\phi(xg)$. Moreover, with this terminology one could say that: ``left operation commutes with right translations''.}

If $\mathscr D\!\mathrm{if}_{G/S}$ denotes the sheaf of $S$-differential operators on $G$ (cf.\ 1.4) and $p$ the structural morphism $G\to S$, one similarly defines a ``right translation'': $\mathscr U(G)\to p_*(\mathscr D\!\mathrm{if}_{G/S})$.

\numpara{2.3}
We shall now characterize the differential operators of $G$ over $S$ of the form $d^G$. Let $g:S\to G$ be a section of the structural morphism of $G$, and $g_G$ the right translation of $G$ by $g$, that is, the composite morphism:
\[
g_G:G\simeq G\times S\xrightarrow{G\times g}G\times G\xrightarrow{m}G.
\]
For every differential operator $D$ of $G$ over $S$, the composite $g_G^{-1}Dg_G$ (cf.\ 1.2) is again an $S$-deviation of $\mathrm{id}_G$, \ie an element of $\mathrm{Dif}_{X/S}$; we shall write:
\[
D^g=g_G^{-1}Dg_G.
\]
% typo?: the source writes Dif_{X/S} here, rather than Dif_{G/S}; retained.
We shall say that $D$ is \emph{right-invariant} if, for every base change $t:T\to S$ and every section $g:T\to G\times T$, one has $(D_T)^g=D_T$.

\begin{enonce}{Lemma}{}
For every differential operator $D$ of $G$ over $S$, the following assertions are equivalent (where $m$ is the multiplication morphism of $G$):
\begin{enumeratei}
\item $D$ is right-invariant.
\item The following two deviations of $m$ are equal: $Dm=m(D\times G)$.
\end{enumeratei}
\end{enonce}

\noindent (ii) $\Rightarrow$ (i): since condition (ii) is stable under base change, it suffices to show that (ii) implies the equality $D^g=D$ for every section $g:S\to G$. Let $h$ be the morphism $G\times g:G\simeq G\times S\to G\times G$, so that $m\circ h$ is the right translation $g_G$. The equality $D^g=D$ is equivalent to the equality $g_G\circ D=D\circ g_G$, and the latter follows from the commutative diagram:
\[
\xymatrix@C=4em@R=3em{G\ar[d]^D_{\mathrm{id}_G}&G\times G\ar[l]_m\ar[d]^{D\times G}_{\mathrm{id}_{(G\times G)}}&G\ar[l]_h\ar[d]^D_{\mathrm{id}_G}\\
G&G\times G\ar[l]^m&G\ar[l]^h.}
\]
\noindent (i) $\Rightarrow$ (ii): take, indeed, for $t:T\to S$ the structural morphism $p:G\to S$, and for the section $g:T\to G\times T$ the diagonal morphism $\Delta:G\to G\times G$. The right translation
\[
\Delta_{G\times G}:G\times G\longrightarrow G\times G
\]
is then the morphism from $G\times G$ to $G\times G$ with components $m$ and $\mathrm{pr}_2$. The equality $(D_G)^\Delta=D_G$ is then equivalent to the commutativity of the first square of the following diagram:
\[
\xymatrix@C=4em@R=3em{G\times G\ar[r]^{\Delta_{G\times G}}\ar[d]^{D_G}_{\mathrm{id}_{G\times G}}&G\times G\ar[r]^{\mathrm{pr}_1}\ar[d]^{D_G}_{\mathrm{id}_{G\times G}}&G\ar[d]^D_{\mathrm{id}_G}\\
G\times G\ar[r]_{\Delta_{G\times G}}&G\times G\ar[r]_{\mathrm{pr}_1}&G.}
\]
Equality (ii) therefore follows from the fact that $m=\mathrm{pr}_1\circ\Delta_{G\times G}$.

Consider, for example, an element $d$ of the infinitesimal algebra $U(G)$. The squares of the diagram
\[
\xymatrix@C=2.2em@R=3em{G\times G\ar@{=}[r]\ar[d]_m&S\times G\times G\ar[r]^{d\times G\times G}_{\varepsilon\times G\times G}\ar[d]_{S\times m}&G\times G\times G\ar[r]^{m\times G}\ar[d]_{G\times m}&G\times G\ar[d]^m\\
G\ar@{=}[r]&S\times G\ar[r]^{d\times G}_{\varepsilon\times G}&G\times G\ar[r]_m&G}
\]
are then commutative. Since one has
\[
m\circ(d\times G)=d^G\qquad\text{and}\qquad(m\times G)\circ(d\times G\times G)=d^G\times G,
\]
one also has $d^G\circ m=m\circ(d^G\times G)$. Thus: \emph{for every $S$-deviation $d$ of the origin, $d^G$ is a right-invariant differential operator}.

\begin{theorem}{2.4}\label{VIIA.2.4}
\begin{enumeratei}
\item The map $d\mapsto d^G$ is an anti-isomorphism\nde{``Isomorphism'' has been corrected to ``anti-isomorphism'', and assertion (ii) has been added, cf.\ N.D.E.\ (17).} from the infinitesimal algebra $U(G)$ onto the subalgebra $\mathrm{Dif}_{G/S}^G$ of $\mathrm{Dif}_{G/S}$ consisting of the right-invariant differential operators.
\item Similarly, the map $d\mapsto{}^Gd$ is an isomorphism from $U(G)$ onto the subalgebra of $\mathrm{Dif}_{G/S}$ consisting of the left-invariant differential operators.
\end{enumeratei}
\end{theorem}

Let, indeed, $D$ be an arbitrary differential operator of $G$ over $S$, and denote by $D_0$ its \emph{value at the origin}, that is, the composite deviation $S\xrightarrow{\varepsilon}G\mathrel{\mathop{\longrightarrow}\limits^D_{\mathrm{id}_G}}G$. The right-invariant differential operator $(D_0)^G$ is then obtained by composition:
\[
\xymatrix@C=2.3em{G\simeq S\times G\ar[r]^{\varepsilon\times G}&G\times G\ar[r]^{D\times G}_{\mathrm{id}_{G\times G}}&G\times G\ar[r]^m&G.}
\]
If $D$ is right-invariant, one has $Dm=m(D\times G)$, whence
\[
D=Dm(\varepsilon\times G)=m(D\times G)(\varepsilon\times G)=(D_0)^G.
\]
In particular, the map $d\mapsto d^G$ is surjective.

Conversely, let $d$ be an $S$-deviation of the origin. There is then a commutative square
\[
\xymatrix@C=4em{G\times G&G\ar[l]_{d\times G}\\
G\times S\simeq G\ar[u]^{G\times\varepsilon}&S\ar[u]_\varepsilon\ar[l]^d}
\]
from which it follows that $d=m(G\times\varepsilon)d=m(d\times G)\varepsilon=(d^G)_0$. A fortiori, the map $d\mapsto d^G$ is injective. This proves the theorem.

As $S$ varies, theorem 2.4 clearly implies that the right translation $\mathscr U(G)\to p_*(\mathscr D\!\mathrm{if}_{G/S})$ is an anti-isomorphism of $\OS$-algebras from $\mathscr U(G)$ onto the sheaf of $\OS$-algebras $p_*(\mathscr D\!\mathrm{if}_{G/S})^G$, which associates $\mathrm{Dif}_{G_U/U}^{G_U}$ to every open subset $U$ of $S$.

\begin{remark}{2.4.1}\label{VIIA.2.4.1}
Consider the commutative diagram
\[
\xymatrix@C=4em@R=3em{G\ar@<0.6ex>[d]_p&G\times G\ar[l]_\eta\ar@<0.6ex>[d]_{\mathrm{pr}_1}\\
S\ar@<0.6ex>[u]_\varepsilon&G\ar[l]^p\ar@<0.6ex>[u]_\Delta,}
\]
where $\eta$ denotes the morphism ``$(x,y)\mapsto yx^{-1}$''.\nde{\ie $G$ acts on the left on itself by right translations.} The latter induces morphisms
\[
\begin{gathered}
\eta':\eta^{-1}(\OO_G)\longrightarrow\OO_{G\times G}\qquad\text{and}\\
\Delta^{-1}(\eta'):p^{-1}\varepsilon^{-1}(\OO_G)\longrightarrow\Delta^{-1}(\OO_{G\times G}).
\end{gathered}
\]
\end{remark}
